\documentclass[12pt,a4paper]{article}

\usepackage[T2A]{fontenc}
\usepackage[utf8]{inputenc}
\usepackage[russian]{babel}
\usepackage{algorithm}
\usepackage{amsmath,amssymb,amsthm,mathtools}
\usepackage{enumitem}
\usepackage{array}
\usepackage{booktabs}
\usepackage{hyperref}
\usepackage{listings}
\usepackage{makecell}
\usepackage{xcolor}

\hypersetup{
    colorlinks=true,
    linkcolor=blue,
    urlcolor=blue
}

\newtheorem{definition}{Определение}
\newtheorem{theorem}{Теорема}
\newtheorem{lemma}{Лемма}
\newtheorem{example}{Пример}

\DeclareMathOperator{\Enc}{Enc}
\DeclareMathOperator{\Dec}{Dec}
\DeclareMathOperator{\Gen}{Gen}
\DeclareMathOperator{\Sign}{Sign}
\DeclareMathOperator{\Ver}{Ver}
\DeclareMathOperator{\Mac}{Mac}
\DeclareMathOperator{\KeyGen}{KeyGen}
\DeclareMathOperator{\Hash}{Hash}
\DeclareMathOperator{\Adv}{Adv}

\title{Билет №18 \\ \small Основные понятия общей алгебры: алгебра, подалгебра, гомоморфизм, конгруэнтность, факторалгебра. Алгебра термов. (СпБПУ)}
\author{Сериков Владислав Александрович}
\date{5 мая 2026}

\begin{document}

\maketitle
\tableofcontents
\newpage

\section{Основные понятия общей алгебры}

\subsection{Сигнатура и алгебра}

\begin{definition}[Сигнатура]
\emph{Сигнатурой} или \emph{типом алгебр} называется множество функциональных символов
\[
\Omega=\{f_i\}_{i\in I},
\]
каждому из которых сопоставлено натуральное число
\[
\operatorname{ar}(f_i)\in \mathbb N_0,
\]
называемое \emph{арностью} символа $f_i$.

Если $\operatorname{ar}(f)=n$, то $f$ называется $n$-арным функциональным символом.
При $n=0$ символ называется \emph{константным}.
\end{definition}

\begin{example}
\begin{enumerate}
    \item Сигнатура групп:
    \[
    \Omega_{\mathrm{Grp}}=\{\cdot,{}^{-1},e\},
    \]
    где
    \[
    \operatorname{ar}(\cdot)=2,\qquad
    \operatorname{ar}({}^{-1})=1,\qquad
    \operatorname{ar}(e)=0.
    \]

    \item Сигнатура колец:
    \[
    \Omega_{\mathrm{Ring}}=\{+,-,\cdot,0,1\},
    \]
    где $+$ и $\cdot$ бинарны, $-$ унарен, $0$ и $1$ являются константами.

    \item Сигнатура решёток:
    \[
    \Omega_{\mathrm{Lat}}=\{\vee,\wedge\},
    \]
    где $\vee$ и $\wedge$ бинарны.
\end{enumerate}
\end{example}

\begin{definition}[Алгебра]
Пусть $\Omega$ --- сигнатура. \emph{Алгеброй сигнатуры $\Omega$} называется пара
\[
\mathcal A=(A,(f^{\mathcal A})_{f\in\Omega}),
\]
где $A$ --- непустое множество, называемое \emph{носителем} или \emph{основным множеством} алгебры, а каждому $n$-арному символу $f\in\Omega$ сопоставлена операция
\[
f^{\mathcal A}\colon A^n\to A.
\]
Если $n=0$, то операция $f^{\mathcal A}\colon A^0\to A$ отождествляется с выделенным элементом множества $A$.
\end{definition}

Обычно, если ясно, о какой алгебре идёт речь, вместо $f^{\mathcal A}$ пишут просто $f$.

\begin{example}
Группа $(G,\cdot,{}^{-1},e)$ является алгеброй сигнатуры
\[
\{\cdot,{}^{-1},e\}.
\]
Кольцо $(R,+,-,\cdot,0,1)$ является алгеброй сигнатуры
\[
\{+,-,\cdot,0,1\}.
\]
\end{example}

\subsection{Подалгебры}

\begin{definition}[Замкнутое подмножество]
Пусть
\[
\mathcal A=(A,(f^{\mathcal A})_{f\in\Omega})
\]
--- алгебра сигнатуры $\Omega$. Подмножество $B\subseteq A$ называется \emph{замкнутым относительно операций алгебры $\mathcal A$}, если для любого $n$-арного символа $f\in\Omega$ и любых элементов
\[
b_1,\dots,b_n\in B
\]
выполнено
\[
f^{\mathcal A}(b_1,\dots,b_n)\in B.
\]
Если $f$ --- константный символ, то условие означает, что
\[
f^{\mathcal A}\in B.
\]
\end{definition}

\begin{definition}[Подалгебра]
Пусть $\mathcal A$ --- алгебра сигнатуры $\Omega$. Алгебра
\[
\mathcal B=(B,(f^{\mathcal B})_{f\in\Omega})
\]
называется \emph{подалгеброй} алгебры $\mathcal A$, если:
\begin{enumerate}
    \item $B\subseteq A$;
    \item $B$ замкнуто относительно всех операций $\mathcal A$;
    \item для каждого $f\in\Omega$ операция $f^{\mathcal B}$ является ограничением операции $f^{\mathcal A}$ на $B$:
    \[
    f^{\mathcal B}=f^{\mathcal A}\big|_{B^n}.
    \]
\end{enumerate}
Обозначение:
\[
\mathcal B\leqslant \mathcal A.
\]
\end{definition}

\begin{example}
Пусть $\mathcal A=(\mathbb Z,+,-,0)$ --- абелева группа целых чисел по сложению. Тогда
\[
2\mathbb Z=\{2k\mid k\in\mathbb Z\}
\]
замкнуто относительно сложения, взятия противоположного элемента и содержит $0$. Поэтому
\[
(2\mathbb Z,+,-,0)
\]
является подалгеброй алгебры $(\mathbb Z,+,-,0)$.
\end{example}

\begin{example}
Множество $\mathbb N\subseteq \mathbb Z$ не является подалгеброй группы $(\mathbb Z,+,-,0)$, поскольку не замкнуто относительно операции $-$:
\[
1\in\mathbb N,\qquad -1\notin\mathbb N.
\]
\end{example}

\subsection{Подалгебра, порождённая множеством}

\begin{definition}[Порождающее множество]
Пусть $\mathcal A$ --- алгебра, $X\subseteq A$. Наименьшая подалгебра алгебры $\mathcal A$, содержащая $X$, называется \emph{подалгеброй, порождённой множеством $X$}, и обозначается
\[
\langle X\rangle_{\mathcal A}.
\]
Если $\langle X\rangle_{\mathcal A}=A$, то говорят, что $X$ \emph{порождает} алгебру $\mathcal A$.
\end{definition}

\begin{theorem}[Существование порождённой подалгебры]
Пусть $\mathcal A$ --- алгебра и $X\subseteq A$. Тогда существует наименьшая подалгебра $\langle X\rangle_{\mathcal A}$, содержащая $X$, причём
\[
\langle X\rangle_{\mathcal A}
=
\bigcap\{B\subseteq A\mid B \text{ является носителем подалгебры } \mathcal A,\ X\subseteq B\}.
\]
\end{theorem}

\begin{proof}
Рассмотрим семейство всех носителей подалгебр алгебры $\mathcal A$, содержащих $X$:
\[
\mathcal S=\{B\subseteq A\mid B \text{ замкнуто относительно всех операций } \mathcal A,\ X\subseteq B\}.
\]
Это семейство непусто, поскольку $A\in\mathcal S$.

Положим
\[
C=\bigcap_{B\in\mathcal S}B.
\]
Очевидно, что $X\subseteq C$, так как $X$ содержится в каждом $B\in\mathcal S$.

Докажем, что $C$ замкнуто относительно всех операций. Пусть $f\in\Omega$ --- $n$-арный символ и
\[
c_1,\dots,c_n\in C.
\]
Тогда для любого $B\in\mathcal S$ имеем
\[
c_1,\dots,c_n\in B.
\]
Так как $B$ замкнуто относительно $f$, то
\[
f^{\mathcal A}(c_1,\dots,c_n)\in B.
\]
Это верно для всех $B\in\mathcal S$, следовательно,
\[
f^{\mathcal A}(c_1,\dots,c_n)\in \bigcap_{B\in\mathcal S}B=C.
\]
Значит, $C$ замкнуто относительно всех операций, то есть является носителем подалгебры.

По построению $C$ содержится в любой подалгебре, содержащей $X$. Следовательно, $C$ является наименьшей такой подалгеброй.
\end{proof}

\subsection{Алгоритм построения подалгебры, порождённой множеством}

Пусть $\mathcal A$ --- конечная алгебра или алгебра, в которой можно эффективно вычислять значения операций, и пусть $X\subseteq A$.

\begin{algorithm}
\caption{Построение $\langle X\rangle_{\mathcal A}$}
\leavevmode
\begin{enumerate}
    \item Положить
    \[
    B_0=X\cup C,
    \]
    где $C$ --- множество всех констант алгебры $\mathcal A$.
    \item Для каждого $n$-арного символа $f\in\Omega$ вычислить все элементы вида
    \[
    f^{\mathcal A}(b_1,\dots,b_n),
    \qquad b_1,\dots,b_n\in B_0.
    \]
    Добавить их к $B_0$. Получится множество $B_1$.
    \item Повторять процесс:
    \[
    B_{k+1}=B_k\cup
    \{f^{\mathcal A}(b_1,\dots,b_n)\mid f\in\Omega,\ b_1,\dots,b_n\in B_k\}.
    \]
    \item Если на некотором шаге
    \[
    B_{k+1}=B_k,
    \]
    то процесс стабилизировался, и
    \[
    \langle X\rangle_{\mathcal A}=B_k.
    \]
\end{enumerate}
\end{algorithm}

\begin{example}
Рассмотрим алгебру
\[
\mathcal A=(\mathbb Z,+,-,0)
\]
и множество
\[
X=\{6,10\}.
\]
Подалгебра, порождённая $6$ и $10$, состоит из всех целочисленных линейных комбинаций этих элементов:
\[
\langle 6,10\rangle_{\mathcal A}
=
\{6m+10n\mid m,n\in\mathbb Z\}.
\]
Так как
\[
\gcd(6,10)=2,
\]
получаем
\[
\langle 6,10\rangle_{\mathcal A}=2\mathbb Z.
\]

Первые шаги:
\[
B_0=\{0,6,10\}.
\]
Замыкая относительно $-$, получаем
\[
-6,\ -10.
\]
Замыкая относительно $+$, получаем
\[
6+10=16,\qquad 10-6=4,\qquad 6-10=-4.
\]
Далее из $6$ и $4$ получаем
\[
6-4=2.
\]
После появления $2$ можно получить все элементы вида $2k$, то есть всю подгруппу $2\mathbb Z$.
\end{example}

\subsection{Гомоморфизмы}

\begin{definition}[Гомоморфизм]
Пусть
\[
\mathcal A=(A,(f^{\mathcal A})_{f\in\Omega}),
\qquad
\mathcal B=(B,(f^{\mathcal B})_{f\in\Omega})
\]
--- алгебры одной и той же сигнатуры $\Omega$.

Отображение
\[
\varphi\colon A\to B
\]
называется \emph{гомоморфизмом} из $\mathcal A$ в $\mathcal B$, если для любого $n$-арного символа $f\in\Omega$ и любых элементов
\[
a_1,\dots,a_n\in A
\]
выполнено
\[
\varphi\bigl(f^{\mathcal A}(a_1,\dots,a_n)\bigr)
=
f^{\mathcal B}\bigl(\varphi(a_1),\dots,\varphi(a_n)\bigr).
\]
Для константного символа $c\in\Omega$ это условие имеет вид
\[
\varphi(c^{\mathcal A})=c^{\mathcal B}.
\]
\end{definition}

\begin{definition}[Мономорфизм, эпиморфизм, изоморфизм]
Гомоморфизм $\varphi\colon \mathcal A\to\mathcal B$ называется:
\begin{enumerate}
    \item \emph{мономорфизмом}, если $\varphi$ инъективен;
    \item \emph{эпиморфизмом}, если $\varphi$ сюръективен;
    \item \emph{изоморфизмом}, если $\varphi$ биективен.
\end{enumerate}

Если существует изоморфизм $\mathcal A\to\mathcal B$, то алгебры $\mathcal A$ и $\mathcal B$ называются \emph{изоморфными}. Обозначение:
\[
\mathcal A\cong\mathcal B.
\]
\end{definition}

\begin{example}
Отображение
\[
\varphi\colon \mathbb Z\to \mathbb Z_n,
\qquad
\varphi(k)=\overline{k}
\]
является гомоморфизмом групп:
\[
(\mathbb Z,+,-,0)\to(\mathbb Z_n,+,-,0).
\]
Действительно,
\[
\varphi(a+b)=\overline{a+b}=\overline a+\overline b=\varphi(a)+\varphi(b),
\]
\[
\varphi(-a)=\overline{-a}=-\overline a=-\varphi(a),
\]
\[
\varphi(0)=\overline 0.
\]
\end{example}

\begin{theorem}[Образ гомоморфизма является подалгеброй]
Пусть
\[
\varphi\colon \mathcal A\to\mathcal B
\]
--- гомоморфизм алгебр сигнатуры $\Omega$. Тогда
\[
\operatorname{Im}\varphi=\{\varphi(a)\mid a\in A\}
\]
является носителем подалгебры алгебры $\mathcal B$.
\end{theorem}

\begin{proof}
Нужно доказать замкнутость $\operatorname{Im}\varphi$ относительно всех операций алгебры $\mathcal B$.

Пусть $f\in\Omega$ --- $n$-арный символ и
\[
b_1,\dots,b_n\in \operatorname{Im}\varphi.
\]
Тогда существуют элементы
\[
a_1,\dots,a_n\in A
\]
такие, что
\[
b_i=\varphi(a_i),\qquad i=1,\dots,n.
\]
По определению гомоморфизма:
\[
f^{\mathcal B}(b_1,\dots,b_n)
=
f^{\mathcal B}(\varphi(a_1),\dots,\varphi(a_n))
=
\varphi(f^{\mathcal A}(a_1,\dots,a_n)).
\]
Последний элемент принадлежит $\operatorname{Im}\varphi$. Значит, образ замкнут относительно $f$.

Для константного символа $c$ имеем
\[
c^{\mathcal B}=\varphi(c^{\mathcal A})\in \operatorname{Im}\varphi.
\]
Следовательно, $\operatorname{Im}\varphi$ является носителем подалгебры алгебры $\mathcal B$.
\end{proof}

\subsection{Конгруэнтности}

\begin{definition}[Отношение эквивалентности]
Бинарное отношение $\theta$ на множестве $A$ называется \emph{отношением эквивалентности}, если для всех $a,b,c\in A$ выполнены свойства:
\begin{enumerate}
    \item рефлексивность:
    \[
    a\mathrel{\theta}a;
    \]
    \item симметричность:
    \[
    a\mathrel{\theta}b \Rightarrow b\mathrel{\theta}a;
    \]
    \item транзитивность:
    \[
    a\mathrel{\theta}b,\ b\mathrel{\theta}c
    \Rightarrow
    a\mathrel{\theta}c.
    \]
\end{enumerate}
\end{definition}

Класс эквивалентности элемента $a\in A$ по отношению $\theta$ обозначается
\[
[a]_\theta=\{x\in A\mid x\mathrel{\theta}a\}.
\]

\begin{definition}[Конгруэнтность]
Пусть
\[
\mathcal A=(A,(f^{\mathcal A})_{f\in\Omega})
\]
--- алгебра сигнатуры $\Omega$.

Отношение эквивалентности $\theta$ на $A$ называется \emph{конгруэнтностью} алгебры $\mathcal A$, если оно согласовано со всеми операциями алгебры, то есть для любого $n$-арного символа $f\in\Omega$ и любых элементов
\[
a_1,\dots,a_n,b_1,\dots,b_n\in A
\]
из условий
\[
a_i\mathrel{\theta}b_i,\qquad i=1,\dots,n,
\]
следует
\[
f^{\mathcal A}(a_1,\dots,a_n)
\mathrel{\theta}
f^{\mathcal A}(b_1,\dots,b_n).
\]
\end{definition}

Множество всех конгруэнтностей алгебры $\mathcal A$ обозначается
\[
\operatorname{Con}(\mathcal A).
\]

\begin{example}
На группе $G$ конгруэнтности находятся во взаимно однозначном соответствии с нормальными подгруппами.

Если $N\trianglelefteq G$, то отношение
\[
a\mathrel{\theta_N}b
\Longleftrightarrow
ab^{-1}\in N
\]
является конгруэнтностью группы $G$.

В частности, для аддитивной группы $\mathbb Z$ и подгруппы $n\mathbb Z$ отношение
\[
a\equiv b\pmod n
\]
является конгруэнтностью.
\end{example}

\begin{example}
Рассмотрим алгебру
\[
(\mathbb Z,+,-,0).
\]
Отношение сравнимости по модулю $3$:
\[
a\mathrel{\theta}b
\Longleftrightarrow
a-b\in 3\mathbb Z
\]
является конгруэнтностью.

Действительно, если
\[
a\equiv b\pmod 3,
\qquad
c\equiv d\pmod 3,
\]
то
\[
a+c\equiv b+d\pmod 3,
\]
а также
\[
-a\equiv -b\pmod 3.
\]
Кроме того, $0\equiv 0\pmod 3$.
\end{example}

\begin{theorem}[Ядро гомоморфизма является конгруэнтностью]
Пусть
\[
\varphi\colon \mathcal A\to\mathcal B
\]
--- гомоморфизм алгебр одной сигнатуры. Определим отношение
\[
\ker\varphi=\{(a,a')\in A^2\mid \varphi(a)=\varphi(a')\}.
\]
Тогда $\ker\varphi$ является конгруэнтностью алгебры $\mathcal A$.
\end{theorem}

\begin{proof}
Сначала докажем, что $\ker\varphi$ является отношением эквивалентности.

\textbf{Рефлексивность.} Для любого $a\in A$ имеем
\[
\varphi(a)=\varphi(a),
\]
следовательно,
\[
(a,a)\in\ker\varphi.
\]

\textbf{Симметричность.} Если
\[
(a,b)\in\ker\varphi,
\]
то
\[
\varphi(a)=\varphi(b).
\]
Следовательно,
\[
\varphi(b)=\varphi(a),
\]
то есть
\[
(b,a)\in\ker\varphi.
\]

\textbf{Транзитивность.} Если
\[
(a,b)\in\ker\varphi
\quad\text{и}\quad
(b,c)\in\ker\varphi,
\]
то
\[
\varphi(a)=\varphi(b),
\qquad
\varphi(b)=\varphi(c).
\]
Значит,
\[
\varphi(a)=\varphi(c),
\]
то есть
\[
(a,c)\in\ker\varphi.
\]

Теперь докажем согласованность с операциями. Пусть $f\in\Omega$ --- $n$-арный символ и
\[
(a_i,b_i)\in\ker\varphi,
\qquad i=1,\dots,n.
\]
Тогда
\[
\varphi(a_i)=\varphi(b_i),
\qquad i=1,\dots,n.
\]
Используя то, что $\varphi$ является гомоморфизмом, получаем:
\[
\varphi(f^{\mathcal A}(a_1,\dots,a_n))
=
f^{\mathcal B}(\varphi(a_1),\dots,\varphi(a_n)).
\]
Так как $\varphi(a_i)=\varphi(b_i)$, то
\[
f^{\mathcal B}(\varphi(a_1),\dots,\varphi(a_n))
=
f^{\mathcal B}(\varphi(b_1),\dots,\varphi(b_n)).
\]
Снова применяя свойство гомоморфизма:
\[
f^{\mathcal B}(\varphi(b_1),\dots,\varphi(b_n))
=
\varphi(f^{\mathcal A}(b_1,\dots,b_n)).
\]
Следовательно,
\[
\varphi(f^{\mathcal A}(a_1,\dots,a_n))
=
\varphi(f^{\mathcal A}(b_1,\dots,b_n)).
\]
Значит,
\[
\bigl(f^{\mathcal A}(a_1,\dots,a_n),
f^{\mathcal A}(b_1,\dots,b_n)\bigr)\in\ker\varphi.
\]
Итак, $\ker\varphi$ является конгруэнтностью.
\end{proof}

\subsection{Факторалгебра}

\begin{definition}[Факторалгебра]
Пусть
\[
\mathcal A=(A,(f^{\mathcal A})_{f\in\Omega})
\]
--- алгебра сигнатуры $\Omega$, а $\theta\in\operatorname{Con}(\mathcal A)$.

\emph{Факторалгеброй} алгебры $\mathcal A$ по конгруэнтности $\theta$ называется алгебра
\[
\mathcal A/\theta
\]
с носителем
\[
A/\theta=\{[a]_\theta\mid a\in A\}
\]
и операциями, заданными правилом
\[
f^{\mathcal A/\theta}([a_1]_\theta,\dots,[a_n]_\theta)
=
[f^{\mathcal A}(a_1,\dots,a_n)]_\theta.
\]
\end{definition}

Необходимо проверить, что это определение корректно: результат не должен зависеть от выбора представителей классов.

\begin{theorem}[Корректность определения операций на факторалгебре]
Пусть $\theta$ --- конгруэнтность алгебры $\mathcal A$. Тогда операции на $A/\theta$, заданные формулой
\[
f^{\mathcal A/\theta}([a_1]_\theta,\dots,[a_n]_\theta)
=
[f^{\mathcal A}(a_1,\dots,a_n)]_\theta,
\]
определены корректно.
\end{theorem}

\begin{proof}
Пусть
\[
[a_i]_\theta=[b_i]_\theta,
\qquad i=1,\dots,n.
\]
Тогда
\[
a_i\mathrel{\theta}b_i,
\qquad i=1,\dots,n.
\]
Так как $\theta$ является конгруэнтностью, то
\[
f^{\mathcal A}(a_1,\dots,a_n)
\mathrel{\theta}
f^{\mathcal A}(b_1,\dots,b_n).
\]
Следовательно,
\[
[f^{\mathcal A}(a_1,\dots,a_n)]_\theta
=
[f^{\mathcal A}(b_1,\dots,b_n)]_\theta.
\]
Значит, значение операции в факторалгебре не зависит от выбора представителей классов эквивалентности.
\end{proof}

\begin{definition}[Естественный гомоморфизм]
Отображение
\[
\pi_\theta\colon A\to A/\theta,
\qquad
\pi_\theta(a)=[a]_\theta,
\]
называется \emph{естественным гомоморфизмом} или \emph{канонической проекцией}.
\end{definition}

\begin{theorem}
Каноническая проекция
\[
\pi_\theta\colon \mathcal A\to \mathcal A/\theta
\]
является сюръективным гомоморфизмом, причём
\[
\ker\pi_\theta=\theta.
\]
\end{theorem}

\begin{proof}
Сюръективность очевидна: каждый элемент факторалгебры имеет вид $[a]_\theta$ для некоторого $a\in A$, и
\[
\pi_\theta(a)=[a]_\theta.
\]

Докажем, что $\pi_\theta$ сохраняет операции. Пусть $f\in\Omega$ --- $n$-арный символ. Тогда
\[
\pi_\theta(f^{\mathcal A}(a_1,\dots,a_n))
=
[f^{\mathcal A}(a_1,\dots,a_n)]_\theta.
\]
С другой стороны,
\[
f^{\mathcal A/\theta}(\pi_\theta(a_1),\dots,\pi_\theta(a_n))
=
f^{\mathcal A/\theta}([a_1]_\theta,\dots,[a_n]_\theta)
=
[f^{\mathcal A}(a_1,\dots,a_n)]_\theta.
\]
Следовательно,
\[
\pi_\theta(f^{\mathcal A}(a_1,\dots,a_n))
=
f^{\mathcal A/\theta}(\pi_\theta(a_1),\dots,\pi_\theta(a_n)).
\]
Значит, $\pi_\theta$ является гомоморфизмом.

Наконец,
\[
(a,b)\in\ker\pi_\theta
\Longleftrightarrow
\pi_\theta(a)=\pi_\theta(b)
\Longleftrightarrow
[a]_\theta=[b]_\theta
\Longleftrightarrow
a\mathrel{\theta}b.
\]
Следовательно,
\[
\ker\pi_\theta=\theta.
\]
\end{proof}

\begin{example}
Рассмотрим группу
\[
(\mathbb Z,+,-,0)
\]
и конгруэнтность
\[
a\equiv b\pmod n.
\]
Тогда факторалгебра
\[
\mathbb Z/\theta
\]
есть обычная группа классов вычетов
\[
\mathbb Z_n.
\]
Операция задаётся правилом
\[
[a]+[b]=[a+b].
\]
\end{example}

\subsection{Первая теорема об изоморфизме}

\begin{theorem}[Первая теорема об изоморфизме]
Пусть
\[
\varphi\colon \mathcal A\to\mathcal B
\]
--- гомоморфизм алгебр одной сигнатуры. Тогда
\[
\mathcal A/\ker\varphi\cong \operatorname{Im}\varphi.
\]
Более точно, отображение
\[
\overline{\varphi}\colon A/\ker\varphi\to \operatorname{Im}\varphi,
\qquad
\overline{\varphi}([a]_{\ker\varphi})=\varphi(a),
\]
является изоморфизмом.
\end{theorem}

\begin{proof}
Обозначим
\[
\theta=\ker\varphi.
\]
Определим отображение
\[
\overline{\varphi}\colon A/\theta\to \operatorname{Im}\varphi
\]
формулой
\[
\overline{\varphi}([a]_\theta)=\varphi(a).
\]

\textbf{1. Корректность.}
Пусть
\[
[a]_\theta=[b]_\theta.
\]
Тогда
\[
a\mathrel{\theta}b.
\]
Так как $\theta=\ker\varphi$, получаем
\[
\varphi(a)=\varphi(b).
\]
Следовательно, значение $\overline{\varphi}([a]_\theta)$ не зависит от выбора представителя класса.

\textbf{2. Сюръективность.}
Пусть
\[
y\in \operatorname{Im}\varphi.
\]
Тогда существует $a\in A$ такое, что
\[
y=\varphi(a).
\]
Но
\[
y=\overline{\varphi}([a]_\theta),
\]
следовательно, $\overline{\varphi}$ сюръективно.

\textbf{3. Инъективность.}
Пусть
\[
\overline{\varphi}([a]_\theta)=\overline{\varphi}([b]_\theta).
\]
Тогда
\[
\varphi(a)=\varphi(b).
\]
Следовательно,
\[
(a,b)\in\ker\varphi=\theta.
\]
Значит,
\[
[a]_\theta=[b]_\theta.
\]
Итак, $\overline{\varphi}$ инъективно.

\textbf{4. Сохранение операций.}
Пусть $f\in\Omega$ --- $n$-арный символ. Тогда
\[
\overline{\varphi}
\left(
f^{\mathcal A/\theta}([a_1]_\theta,\dots,[a_n]_\theta)
\right)
=
\overline{\varphi}
\left(
[f^{\mathcal A}(a_1,\dots,a_n)]_\theta
\right).
\]
По определению $\overline{\varphi}$:
\[
\overline{\varphi}
\left(
[f^{\mathcal A}(a_1,\dots,a_n)]_\theta
\right)
=
\varphi(f^{\mathcal A}(a_1,\dots,a_n)).
\]
Так как $\varphi$ --- гомоморфизм,
\[
\varphi(f^{\mathcal A}(a_1,\dots,a_n))
=
f^{\mathcal B}(\varphi(a_1),\dots,\varphi(a_n)).
\]
Но
\[
\varphi(a_i)=\overline{\varphi}([a_i]_\theta),
\]
поэтому
\[
f^{\mathcal B}(\varphi(a_1),\dots,\varphi(a_n))
=
f^{\operatorname{Im}\varphi}
(
\overline{\varphi}([a_1]_\theta),
\dots,
\overline{\varphi}([a_n]_\theta)
).
\]
Следовательно, $\overline{\varphi}$ является гомоморфизмом.

Так как $\overline{\varphi}$ является биективным гомоморфизмом, оно является изоморфизмом.
\end{proof}

\begin{example}
Рассмотрим гомоморфизм
\[
\varphi\colon \mathbb Z\to \mathbb Z_n,
\qquad
\varphi(k)=\overline{k}.
\]
Тогда
\[
\ker\varphi=n\mathbb Z.
\]
По первой теореме об изоморфизме:
\[
\mathbb Z/n\mathbb Z\cong \mathbb Z_n.
\]
\end{example}

\subsection{Термы}

\begin{definition}[Множество переменных]
Пусть
\[
X=\{x_1,x_2,x_3,\dots\}
\]
--- множество переменных. Переменные используются как формальные символы, из которых с помощью операций сигнатуры строятся термы.
\end{definition}

\begin{definition}[Терм]
Пусть $\Omega$ --- сигнатура, а $X$ --- множество переменных. Множество термов $T_\Omega(X)$ определяется индуктивно:
\begin{enumerate}
    \item всякая переменная $x\in X$ является термом;
    \item если $c\in\Omega$ --- константный символ, то $c$ является термом;
    \item если $f\in\Omega$ --- $n$-арный функциональный символ и $t_1,\dots,t_n$ --- термы, то выражение
    \[
    f(t_1,\dots,t_n)
    \]
    является термом;
    \item других термов нет.
\end{enumerate}
\end{definition}

\begin{example}
Пусть сигнатура групп:
\[
\Omega=\{\cdot,{}^{-1},e\}.
\]
Тогда выражения
\[
x,\qquad e,\qquad x^{-1},\qquad x\cdot y,\qquad (x\cdot y^{-1})\cdot z
\]
являются термами.
\end{example}

\begin{example}
Пусть сигнатура колец:
\[
\Omega=\{+,-,\cdot,0,1\}.
\]
Тогда термами являются:
\[
x+y,\qquad x\cdot y+1,\qquad -(x\cdot x+y).
\]
Формально, например, выражение
\[
x\cdot y+1
\]
записывается как
\[
+( \cdot(x,y),1).
\]
\end{example}

\subsection{Интерпретация термов в алгебре}

Пусть
\[
\mathcal A=(A,(f^{\mathcal A})_{f\in\Omega})
\]
--- алгебра сигнатуры $\Omega$.

Каждому терму
\[
t=t(x_1,\dots,x_n)
\]
соответствует операция на $A$, называемая \emph{термальной операцией}:
\[
t^{\mathcal A}\colon A^n\to A.
\]

Она определяется индукцией по построению терма:
\begin{enumerate}
    \item если $t=x_i$, то
    \[
    t^{\mathcal A}(a_1,\dots,a_n)=a_i;
    \]
    \item если $t=c$ --- константа, то
    \[
    t^{\mathcal A}(a_1,\dots,a_n)=c^{\mathcal A};
    \]
    \item если
    \[
    t=f(t_1,\dots,t_m),
    \]
    то
    \[
    t^{\mathcal A}(a_1,\dots,a_n)
    =
    f^{\mathcal A}
    \bigl(
    t_1^{\mathcal A}(a_1,\dots,a_n),
    \dots,
    t_m^{\mathcal A}(a_1,\dots,a_n)
    \bigr).
    \]
\end{enumerate}

\begin{example}
В группе $G$ терм
\[
t(x,y)=xy^{-1}
\]
задаёт бинарную операцию
\[
t^G(a,b)=ab^{-1}.
\]
\end{example}

\subsection{Алгебра термов}

\begin{definition}[Алгебра термов]
Пусть $\Omega$ --- сигнатура, а $X$ --- множество переменных. \emph{Алгеброй термов} называется алгебра
\[
\mathbf T_\Omega(X)
=
(T_\Omega(X),(f^{\mathbf T})_{f\in\Omega}),
\]
где носителем является множество всех термов $T_\Omega(X)$, а операции определяются формально:
если $f\in\Omega$ имеет арность $n$, то
\[
f^{\mathbf T}(t_1,\dots,t_n)=f(t_1,\dots,t_n).
\]
\end{definition}

Иными словами, операции в алгебре термов не вычисляют значения, а строят новые формальные выражения.

\begin{example}
Пусть
\[
\Omega=\{+,-,0\},
\]
где $+$ бинарен, $-$ унарен, $0$ --- константа. Тогда в алгебре термов:
\[
+^{\mathbf T}(x,y)=+(x,y),
\]
\[
-^{\mathbf T}(+(x,y))=-(+(x,y)).
\]
В привычной записи это:
\[
x+y,
\qquad
-(x+y).
\]
\end{example}

\subsection{Свободность алгебры термов}

Алгебра термов является фундаментальным примером свободной алгебры.

\begin{definition}[Свободная алгебра]
Пусть $\mathcal K$ --- класс алгебр сигнатуры $\Omega$. Алгебра $\mathcal F\in\mathcal K$ называется \emph{свободной в классе $\mathcal K$ с множеством свободных порождающих $X$}, если задано вложение
\[
\iota\colon X\to F
\]
и для любой алгебры $\mathcal A\in\mathcal K$ и любого отображения
\[
\alpha\colon X\to A
\]
существует единственный гомоморфизм
\[
\widehat{\alpha}\colon \mathcal F\to \mathcal A
\]
такой, что
\[
\widehat{\alpha}\circ\iota=\alpha.
\]
\end{definition}

Для класса всех алгебр данной сигнатуры свободной алгеброй является алгебра термов.

\begin{theorem}[Универсальное свойство алгебры термов]
Пусть $\Omega$ --- сигнатура, $X$ --- множество переменных, а
\[
\mathbf T_\Omega(X)
\]
--- алгебра термов. Тогда для любой $\Omega$-алгебры $\mathcal A$ и любого отображения
\[
\alpha\colon X\to A
\]
существует единственный гомоморфизм
\[
\widehat{\alpha}\colon \mathbf T_\Omega(X)\to \mathcal A
\]
такой, что
\[
\widehat{\alpha}(x)=\alpha(x)
\]
для всех $x\in X$.
\end{theorem}

\begin{proof}
Построим отображение $\widehat{\alpha}$ индукцией по построению терма.

\textbf{1. Определение отображения.}

Если $t=x\in X$, то полагаем
\[
\widehat{\alpha}(x)=\alpha(x).
\]

Если $t=c$, где $c$ --- константный символ, то полагаем
\[
\widehat{\alpha}(c)=c^{\mathcal A}.
\]

Если
\[
t=f(t_1,\dots,t_n),
\]
то полагаем
\[
\widehat{\alpha}(t)
=
f^{\mathcal A}
\bigl(
\widehat{\alpha}(t_1),\dots,\widehat{\alpha}(t_n)
\bigr).
\]

Таким образом, $\widehat{\alpha}$ определено на всех термах.

\textbf{2. Проверка, что $\widehat{\alpha}$ является гомоморфизмом.}

Пусть $f\in\Omega$ --- $n$-арный символ, и пусть
\[
t_1,\dots,t_n\in T_\Omega(X).
\]
В алгебре термов
\[
f^{\mathbf T}(t_1,\dots,t_n)=f(t_1,\dots,t_n).
\]
Тогда по определению $\widehat{\alpha}$:
\[
\widehat{\alpha}
\bigl(
f^{\mathbf T}(t_1,\dots,t_n)
\bigr)
=
\widehat{\alpha}(f(t_1,\dots,t_n))
=
f^{\mathcal A}
\bigl(
\widehat{\alpha}(t_1),\dots,\widehat{\alpha}(t_n)
\bigr).
\]
Это и означает, что $\widehat{\alpha}$ сохраняет операцию $f$.

\textbf{3. Согласование с $\alpha$.}

Для каждой переменной $x\in X$ по построению:
\[
\widehat{\alpha}(x)=\alpha(x).
\]

\textbf{4. Единственность.}

Пусть
\[
\psi\colon \mathbf T_\Omega(X)\to \mathcal A
\]
--- другой гомоморфизм такой, что
\[
\psi(x)=\alpha(x)
\]
для всех $x\in X$.

Докажем индукцией по построению терма, что
\[
\psi(t)=\widehat{\alpha}(t)
\]
для любого терма $t$.

Если $t=x\in X$, то
\[
\psi(x)=\alpha(x)=\widehat{\alpha}(x).
\]

Если $t=c$ --- константа, то, поскольку $\psi$ --- гомоморфизм,
\[
\psi(c)=c^{\mathcal A}=\widehat{\alpha}(c).
\]

Пусть
\[
t=f(t_1,\dots,t_n),
\]
и для термов $t_1,\dots,t_n$ уже доказано, что
\[
\psi(t_i)=\widehat{\alpha}(t_i),
\qquad i=1,\dots,n.
\]
Тогда
\[
\psi(t)
=
\psi(f^{\mathbf T}(t_1,\dots,t_n)).
\]
Так как $\psi$ --- гомоморфизм,
\[
\psi(f^{\mathbf T}(t_1,\dots,t_n))
=
f^{\mathcal A}(\psi(t_1),\dots,\psi(t_n)).
\]
По предположению индукции:
\[
f^{\mathcal A}(\psi(t_1),\dots,\psi(t_n))
=
f^{\mathcal A}(\widehat{\alpha}(t_1),\dots,\widehat{\alpha}(t_n)).
\]
По определению $\widehat{\alpha}$:
\[
f^{\mathcal A}(\widehat{\alpha}(t_1),\dots,\widehat{\alpha}(t_n))
=
\widehat{\alpha}(f(t_1,\dots,t_n)).
\]
Следовательно,
\[
\psi(t)=\widehat{\alpha}(t).
\]
Значит, $\psi=\widehat{\alpha}$.

Итак, требуемый гомоморфизм существует и единственен.
\end{proof}

\subsection{Связь термов, гомоморфизмов и конгруэнтностей}

\begin{theorem}[Конгруэнтность сохраняется термами]
Пусть $\theta$ --- конгруэнтность алгебры $\mathcal A$. Если
\[
a_i\mathrel{\theta}b_i,
\qquad i=1,\dots,n,
\]
то для любого терма
\[
t(x_1,\dots,x_n)
\]
выполнено
\[
t^{\mathcal A}(a_1,\dots,a_n)
\mathrel{\theta}
t^{\mathcal A}(b_1,\dots,b_n).
\]
\end{theorem}

\begin{proof}
Докажем утверждение индукцией по построению терма $t$.

\textbf{База индукции.}

Если $t=x_i$, то
\[
t^{\mathcal A}(a_1,\dots,a_n)=a_i,
\qquad
t^{\mathcal A}(b_1,\dots,b_n)=b_i.
\]
По условию
\[
a_i\mathrel{\theta}b_i.
\]

Если $t=c$ --- константа, то
\[
t^{\mathcal A}(a_1,\dots,a_n)=c^{\mathcal A}
=
t^{\mathcal A}(b_1,\dots,b_n).
\]
Следовательно,
\[
c^{\mathcal A}\mathrel{\theta}c^{\mathcal A}
\]
по рефлексивности $\theta$.

\textbf{Переход индукции.}

Пусть
\[
t=f(t_1,\dots,t_m).
\]
Предположим, что для термов $t_1,\dots,t_m$ утверждение уже доказано:
\[
t_j^{\mathcal A}(a_1,\dots,a_n)
\mathrel{\theta}
t_j^{\mathcal A}(b_1,\dots,b_n),
\qquad j=1,\dots,m.
\]
Так как $\theta$ является конгруэнтностью, то из этих соотношений следует
\[
f^{\mathcal A}
\bigl(
t_1^{\mathcal A}(a_1,\dots,a_n),
\dots,
t_m^{\mathcal A}(a_1,\dots,a_n)
\bigr)
\mathrel{\theta}
f^{\mathcal A}
\bigl(
t_1^{\mathcal A}(b_1,\dots,b_n),
\dots,
t_m^{\mathcal A}(b_1,\dots,b_n)
\bigr).
\]
Но левая часть есть
\[
t^{\mathcal A}(a_1,\dots,a_n),
\]
а правая часть есть
\[
t^{\mathcal A}(b_1,\dots,b_n).
\]
Следовательно,
\[
t^{\mathcal A}(a_1,\dots,a_n)
\mathrel{\theta}
t^{\mathcal A}(b_1,\dots,b_n).
\]
Теорема доказана.
\end{proof}

\subsection{Алгоритм проверки, является ли отношение конгруэнтностью}

Пусть дана конечная алгебра
\[
\mathcal A=(A,(f^{\mathcal A})_{f\in\Omega})
\]
и бинарное отношение $\theta\subseteq A^2$.

\begin{algorithm}
\caption{Проверка конгруэнтности}
\leavevmode
\begin{enumerate}
    \item Проверить, что $\theta$ является отношением эквивалентности:
    \begin{enumerate}
        \item для всех $a\in A$ проверить $(a,a)\in\theta$;
        \item для всех $a,b\in A$ проверить:
        \[
        (a,b)\in\theta\Rightarrow (b,a)\in\theta;
        \]
        \item для всех $a,b,c\in A$ проверить:
        \[
        (a,b)\in\theta,\ (b,c)\in\theta
        \Rightarrow
        (a,c)\in\theta.
        \]
    \end{enumerate}

    \item Для каждой операции $f\in\Omega$ арности $n$ проверить совместимость:
    для всех наборов
    \[
    (a_1,\dots,a_n),\quad (b_1,\dots,b_n)\in A^n
    \]
    таких, что
    \[
    (a_i,b_i)\in\theta,\qquad i=1,\dots,n,
    \]
    проверить:
    \[
    \bigl(
    f^{\mathcal A}(a_1,\dots,a_n),
    f^{\mathcal A}(b_1,\dots,b_n)
    \bigr)\in\theta.
    \]

    \item Если обе проверки успешны, то $\theta$ является конгруэнтностью. Иначе не является.
\end{enumerate}
\end{algorithm}

\begin{example}
Рассмотрим алгебру
\[
\mathcal A=(\mathbb Z_4,+,-,0).
\]
Пусть отношение $\theta$ разбивает $\mathbb Z_4$ на классы
\[
\{0,2\},\qquad \{1,3\}.
\]
Это отношение соответствует сравнению по модулю $2$:
\[
a\mathrel{\theta}b
\Longleftrightarrow
a\equiv b\pmod 2.
\]

Проверим совместимость со сложением. Если
\[
a\equiv b\pmod 2,
\qquad
c\equiv d\pmod 2,
\]
то
\[
a+c\equiv b+d\pmod 2.
\]
Значит,
\[
(a+c)\mathrel{\theta}(b+d).
\]

Проверим совместимость с унарной операцией $-$. Если
\[
a\equiv b\pmod 2,
\]
то
\[
-a\equiv -b\pmod 2.
\]
Константа $0$ автоматически совместима.

Следовательно, $\theta$ является конгруэнтностью.
\end{example}

\subsection{Минимальная и максимальная конгруэнтности}

Для любой алгебры $\mathcal A$ существуют две очевидные конгруэнтности.

\begin{definition}
\emph{Диагональная конгруэнтность}:
\[
\Delta_A=\{(a,a)\mid a\in A\}.
\]
Она отождествляет только равные элементы.

\emph{Универсальная конгруэнтность}:
\[
\nabla_A=A\times A.
\]
Она отождествляет все элементы алгебры.
\end{definition}

Очевидно,
\[
\Delta_A\subseteq \theta\subseteq \nabla_A
\]
для любой конгруэнтности $\theta\in\operatorname{Con}(\mathcal A)$.

Факторалгебры по этим конгруэнтностям имеют вид:
\[
\mathcal A/\Delta_A\cong \mathcal A,
\]
\[
\mathcal A/\nabla_A
\]
--- одноэлементная алгебра.

\subsection{Конгруэнтность, порождённая отношением}

\begin{definition}
Пусть $R\subseteq A^2$. \emph{Конгруэнтностью, порождённой отношением $R$}, называется наименьшая конгруэнтность $\theta$ алгебры $\mathcal A$, такая что
\[
R\subseteq \theta.
\]
Она обозначается
\[
\operatorname{Cg}^{\mathcal A}(R).
\]
\end{definition}

\begin{theorem}[Существование порождённой конгруэнтности]
Для любого отношения $R\subseteq A^2$ существует наименьшая конгруэнтность, содержащая $R$, причём
\[
\operatorname{Cg}^{\mathcal A}(R)
=
\bigcap\{\theta\in\operatorname{Con}(\mathcal A)\mid R\subseteq\theta\}.
\]
\end{theorem}

\begin{proof}
Семейство
\[
\{\theta\in\operatorname{Con}(\mathcal A)\mid R\subseteq\theta\}
\]
непусто, так как универсальная конгруэнтность $\nabla_A=A^2$ содержит $R$.

Положим
\[
\Theta=
\bigcap\{\theta\in\operatorname{Con}(\mathcal A)\mid R\subseteq\theta\}.
\]

Пересечение отношений эквивалентности снова является отношением эквивалентности. Проверим совместимость с операциями.

Пусть $f\in\Omega$ --- $n$-арная операция, и пусть
\[
(a_i,b_i)\in\Theta,
\qquad i=1,\dots,n.
\]
Это означает, что для любой конгруэнтности $\theta$, содержащей $R$, выполнено
\[
(a_i,b_i)\in\theta.
\]
Так как каждая такая $\theta$ является конгруэнтностью, получаем
\[
\bigl(
f^{\mathcal A}(a_1,\dots,a_n),
f^{\mathcal A}(b_1,\dots,b_n)
\bigr)\in\theta
\]
для всех $\theta$, содержащих $R$. Следовательно,
\[
\bigl(
f^{\mathcal A}(a_1,\dots,a_n),
f^{\mathcal A}(b_1,\dots,b_n)
\bigr)\in\Theta.
\]
Значит, $\Theta$ является конгруэнтностью.

Она содержит $R$, поскольку каждое отношение из пересекаемого семейства содержит $R$. Кроме того, по построению $\Theta$ содержится в любой конгруэнтности, содержащей $R$. Следовательно, $\Theta$ является наименьшей такой конгруэнтностью.
\end{proof}

\subsection{Итоговая схема основных понятий}

Основные конструкции общей алгебры связаны следующим образом:
\[
\boxed{
\text{сигнатура}
\Rightarrow
\text{алгебра}
\Rightarrow
\text{подалгебры, гомоморфизмы, конгруэнтности}
}
\]

\[
\boxed{
\text{конгруэнтность}
\Rightarrow
\text{факторалгебра}
}
\]

\[
\boxed{
\text{гомоморфизм}
\Rightarrow
\text{ядро}
\Rightarrow
\text{конгруэнтность}
}
\]

\[
\boxed{
\mathcal A/\ker\varphi\cong\operatorname{Im}\varphi
}
\]

\[
\boxed{
\text{термы}
\Rightarrow
\text{алгебра термов}
\Rightarrow
\text{свободная алгебра}
}
\]

Таким образом, общая алгебра изучает алгебраические структуры с произвольным набором операций, не завися от конкретной природы этих операций. Центральными объектами являются алгебры фиксированной сигнатуры, их подалгебры, гомоморфизмы, конгруэнтности и факторалгебры. Алгебра термов играет роль свободной алгебры и позволяет формально описывать все выражения, которые можно построить из переменных и операций данной сигнатуры.


\end{document}
